\documentclass{article}

\usepackage{amsmath}
\usepackage{amsthm}
\usepackage{amssymb}
\usepackage{hyperref}
\usepackage{listings}
\usepackage{xcolor}

\lstset{
    language=Python,
    basicstyle=\ttfamily\footnotesize,
    keywordstyle=\color{blue},
    commentstyle=\color{green!60!black},
    stringstyle=\color{red},
    showstringspaces=false,
    breaklines=true,
    frame=single,
    numbers=left,
    numberstyle=\tiny\color{gray},
    tabsize=4,
    columns=flexible,
    keepspaces=true,
    basewidth=0.5em
}

\usepackage{enumitem}
\setlist[enumerate,1]{label=\arabic*.}
\setlist[enumerate,2]{label=(\alph*)}
\setlist[enumerate,3]{label=\roman*.}

\title{Written HW3}
\author{
    Logan Tanner \\ 
    I use Neovim (BTW) \\ 
    \\
    Full HW repository at \\
    \href{https://github.com/LoganCTanner/School/tree/main/Math/425-Probability}{github.com/LoganCTanner/School}
}
\date{May 15, 2026}

\begin{document}
\maketitle

\clearpage
\begin{enumerate}

    \item   Suppose the number of times a person contracts a cold in a given year is a Poisson with $\lambda$ = 5.
            Also Suppose a new wonder drug has just been marketed that reduces the Poisson parameter to $\lambda$ = 3
            for 75\% of the population; for the other 25\%, the drug has no effect. If an individual tries the drug 
            for a year, and gets 2 colds during that time, what is the probability that the drug is beneficial for 
            this individual?

            
            Answer:\\
            To find the probability that the drug is beneficial for this individual we may use Bayes' Theorem to find
            the likelihood that that they are in the group of individuals where the 
            poisson distribution is given by $\lambda$ = 3.\\
            $ P(\lambda = 3 | \text{drug}) = .75\text{, } P(\lambda = 5 | \text{drug}) = .25$\\
            $ P(\lambda = 3 | k = 2) = \frac{P(\lambda = 3)P(k=2 | \lambda=3)}{P(k = 2)} $
            $$ P(k = 2) = P(\lambda = 3)P(k=2 | \lambda=3) + P(\lambda=5)P(k = 2 | \lambda=5) $$
            $$ = .75\left(e^{-3}\frac{3^2}{2!}\right) + .25\left(e^{-5}\frac{5^2}{2!}\right) $$
            $$ \approx .75\times.224042 + .25\times.0842243 $$
            $$ \approx .189088 $$
            $$ P(\lambda = 3 | k = 2) \approx \frac{.75 \times .224042}{.189088} $$
            $$ \approx \fbox{.888641} $$

    \clearpage
    \item   An Interviewer is given a list of people she can interview. She will conduct at most 5 interviews, 
            and each person independently agrees to be interviewed with probability $\frac{2}{3}$.
            
            \begin{enumerate}
                \item   If the list has 5 people on it, what's the probability that she gets five interviews?
                        \\Answer: \\
                        This question follows the format of a Binomial Distribution Probability for the simplest case 
                        where $k = n$ \\
                        $ X \sim Binomial(5,\frac{2}{3}) $\\
                        $ P(X=5) = \binom{5}{5} \times \left(\frac{2}{3}\right)^{5} \times \left(\frac{1}{3}\right)^{0} $
                        $$ = \framebox{ $\frac{32}{243}$} $$
                \item   What if the list has 8 people on it?
                        \\Answer: \\
                        The question remains to be a binomial, though now with $ 0 < k < n$, resulting in\\
                        $ X \sim Binomial(8,\frac{2}{3}) $\\
                        $ P(X=5) = \binom{8}{5} \times \left(\frac{2}{3}\right)^{5} \times \left(\frac{1}{3}\right)^{3} $
                        $$ = \framebox{ $\frac{1792}{6561}$} $$
                \clearpage
                \item   If the list has 8 people on it, and she begins at the top, 
                        what is the probability that she interviews the 7th person?
                        \\Answer:\\
                        In this case there are 6 possible people preceding the 7th interviewee 
                        $ X \sim Binomial(6,\frac{2}{3}) $
                        $$ P(\text{7th person interviewed}) = \frac{2}{3} \times \Sigma_{i=0}^{4}P(X=i) $$
                        $$ P(X=i) = \binom{6}{i} \times \left(\frac{2}{3}\right)^i \times \left(\frac{1}{3}\right)^{6-i} $$
                        $$  \Sigma_{i=0}^{4} P(X=i) = $$ 
                        $$ \binom{6}{0} \times \left(\frac{2}{3}\right)^0 \times \left(\frac{1}{3}\right)^6 $$
                        $$  + \binom{6}{1} \times \left(\frac{2}{3}\right)^1 \times \left(\frac{1}{3}\right)^5 $$
                        $$  + \binom{6}{2} \times \left(\frac{2}{3}\right)^2 \times \left(\frac{1}{3}\right)^4 $$
                        $$  + \binom{6}{3} \times \left(\frac{2}{3}\right)^3 \times \left(\frac{1}{3}\right)^3 $$
                        $$  + \binom{6}{4} \times \left(\frac{2}{3}\right)^4 \times \left(\frac{1}{3}\right)^2 $$
                        $$ = \frac{1}{3^6} + \frac{6\times2}{3^6} + \frac{15\times2^2}{3^6} + \frac{20 \times 2^3}{3^6} + \frac{15\times 2^4}{3^6} $$
                        $$ = \frac{473}{729} $$
                        $$ P(\text{7th person interviewed}) = \frac{2}{3} \times \frac{473}{729} = \fbox{$ \frac{946}{2187} $} \approx 0.432556 $$
                        To verify, I ran the following Python script, which returned $0.433336$
                        \clearpage
                        \begin{lstlisting}
import random

count = 0
trials = 1000000

for _ in range(trials):
    interviews = 0
    
    # Go through persons 1-6
    for person in range(1, 7):
        if random.random() < 2/3:  # Person says yes
            interviews += 1
            if interviews == 5:  # Hit the limit
                break
    
    # Check if we can reach person 7
    if interviews < 5:
        # Ask person 7
        if random.random() < 2/3:  # Person 7 says yes
            count += 1

probability = count / trials
print(f"P(interviews person 7) = {probability}")
                        \end{lstlisting}
                \clearpage
                \item   As in part (c), what is the probability she interviews the 8th person?
                    \\Answer:\\
                    Like in (c), we are looking for the probability where the $j^{\text{th}}$ positioned
                    potential interviewee is interviewed, where now $j = 8$.
                    We can then substitute in $j - 1= 7$ for the number of preceding potential interviewees
                    and recalculate to find:\\
                    $ X \sim Binomial(7,\frac{2}{3}) $\\
                    $ P(\text{8th person interviewed}) = \frac{2}{3} \times \Sigma_{i=0}^{4}P(X=i) $\\
                    $ P(X=i) = \binom{7}{i} \times \left(\frac{2}{3}\right)^i \times \left(\frac{1}{3}\right)^{7-i} $\\
                    $ \Rightarrow P(\text{8th person interviewed}) = \fbox{$\frac{1878}{6561}$} \approx 0.28624 $ \\
                    Then, modifying the previous Python script to be range(1,8) on line 10 and recomputing, 
                    the approximation verifies with a result of 0.285745


    
            \end{enumerate}


\end{enumerate}

\end{document}
